\section{Do quality and content filters have similar effects?}
\label{sec:quality-content-correlation}

In order to further understand how filters described in \autoref{sub:web:quality}, \autoref{sub:web:content}, and \autoref{sub:web:dedup} interact with each other, we perform a correlation analysis on a subset of documents sampled from our pipeline. 
The correlation among the documents flagged for removal by our Common Crawl filters is depicted in Figure~\ref{fig:corr}.
Overall, we find that correlations are generally low, thus our filters select fairly different documents and are not redundant. 


\begin{figure*}[htbp]
    \centering
    \subfloat[\centering High]{{\includegraphics[width=0.25\textwidth]{figures/heatmap_head.pdf}}}
    \qquad
    {\subfloat[\centering Medium]{{\includegraphics[width=0.25\textwidth]{figures/heatmap_middle.pdf}}}}
    \qquad
    {\subfloat[\centering Low]{{\includegraphics[width=0.25\textwidth]{figures/heatmap_tail.pdf}}}}    
    \caption{Pearson Correlation of various \dolma filters on the High, Medium, and Low buckets of our Common Crawl data, computed over 24M, 20M, and 43M documents, respectively. The filters are Gopher=Gopher rules from \citet{Rae2021ScalingLM}, Dedup.=Deduplication, PII=Personally Identifiable Information, Hate=Toxicity and Decont.=Decontamination.
    Calculated at the document-level: two filters contribute to positive correlation when any span in a document is tagged by both filters.
    We find our various filters remove different documents and are not redundant.
    }
    \label{fig:corr}
\end{figure*}


There is some positive correlation between our PII (Personal Identifiable Information) filters and filters removing hate speech. 
This is likely because hate speech is often directed at people. 
The Gopher filtering rules correlate negatively with our deduplication, especially for the high-perplexity tail part of our data. 
This is due to the Gopher rules removing many high-perplexity documents such as random strings, which are not caught by deduplication due to their randomness. 
As these random strings likely do not contribute to a better understanding of language, it is important to filter them out and thus rely on filters beyond deduplication. 
